% =============================================================================
% tema/tikz.tex - vocabulario grafico da serie
% -----------------------------------------------------------------------------
% Conjunto musial*, herdado sem ruptura da implementacao de referencia
% (FIBROMIALGIA_2026_REFINADA): gradiente branco para tonalidade clara, borda
% semantica, sombra difusa discreta e seta com a cor do destino.
%
% Acrescimos validados na serie e mantidos aqui: supressao de hifenizacao em
% todo no, os atalhos crnode, crseta e passo, a moldura \bandaimg e o \chip.
%
% Regra que vale em todo frame: a seta sai de uma FACE explicita do no
% (.south, .east, .north, .west) e nunca do centro. Uma seta ancorada no centro
% calcula o ++ a partir do centro e, quando a altura minima do no cresce, o
% rotulo cai dentro da propria caixa.
% =============================================================================

\tikzset{
  % --- base -----------------------------------------------------------------
  musialBase/.style={rectangle, rounded corners=3pt, very thick, align=center,
                     inner xsep=6pt, inner ysep=5pt, font=\footnotesize,
                     execute at begin node={\hyphenpenalty=10000\exhyphenpenalty=10000},
                     blur shadow={shadow blur steps=4, shadow xshift=0.8pt,
                                  shadow yshift=-0.8pt}},
  musialBlue/.style={musialBase, draw=campoazul!80!black,
                     top color=white, bottom color=campoazul!18},
  musialGreen/.style={musialBase, draw=campoverde!82!black,
                      top color=white, bottom color=campoverde!18},
  musialRed/.style={musialBase, draw=campobordo!88!black,
                    top color=white, bottom color=campobordo!14},
  musialAmber/.style={musialBase, draw=campoambar!88!black,
                      top color=white, bottom color=campoambar!18},
  musialPurple/.style={musialBase, draw=camporoxo!80!black,
                       top color=white, bottom color=camporoxo!16},
  musialYellow/.style={musialBase, draw=campoamarelo!72!black,
                       top color=white, bottom color=campoamarelo!24},
  musialNeutral/.style={musialBase, draw=campografite!55,
                        top color=white, bottom color=campografite!6},
  % --- circulo central ------------------------------------------------------
  musialCenter/.style={circle, very thick, draw=campoazul!85!black,
                       top color=white, bottom color=campoazul!28,
                       text=campoazul!90!black, align=center,
                       font=\small\bfseries\sffamily,
                       minimum size=2.25cm, inner sep=4pt,
                       execute at begin node={\hyphenpenalty=10000\exhyphenpenalty=10000},
                       blur shadow={shadow blur steps=5, shadow xshift=1pt,
                                    shadow yshift=-1pt}},
  % --- setas ----------------------------------------------------------------
  musialFlow/.style={-{Latex[length=2.5mm,width=1.8mm]}, line width=0.8pt,
                     campoazul!85!black, rounded corners=4pt},
  musialFlowGreen/.style={musialFlow, campoverde!85!black},
  musialFlowRed/.style={musialFlow, campobordo!88!black},
  musialFlowAmber/.style={musialFlow, campoambar!90!black},
  musialFlowPurple/.style={musialFlow, camporoxo!85!black},
  musialFlowGraf/.style={musialFlow, campografite!80},
  musialThin/.style={-{Latex[length=2.0mm]}, semithick, campografite!70},
  musialBi/.style={{Latex[length=2.2mm,width=1.6mm]}-{Latex[length=2.2mm,width=1.6mm]},
                   line width=0.7pt, campoazul!75!black, rounded corners=4pt},
  % --- dose, herdados -------------------------------------------------------
  musialDose/.style={musialBlue, text width=3.0cm, minimum height=0.95cm,
                     font=\footnotesize},
  musialDoseTarget/.style={musialGreen, text width=3.0cm, minimum height=0.95cm,
                           font=\footnotesize},
  % --- rotulo sobre linha ---------------------------------------------------
  musialRot/.style={font=\scriptsize\sffamily, text=campografite,
                    fill=white, fill opacity=0.9, text opacity=1,
                    inner sep=1.6pt, rounded corners=1pt, align=center},
  % --- chip inline ----------------------------------------------------------
  chipno/.style args={#1}{rectangle, rounded corners=2pt, draw=#1!75!black,
                          fill=#1!10, text=#1!80!black, inner xsep=3.2pt,
                          inner ysep=1.2pt, font=\scriptsize\sffamily\bfseries},
}

% =============================================================================
% ATALHOS DE AUTOR
% -----------------------------------------------------------------------------
% crnode={cor}{largura} : no da familia musial na cor semantica pedida, com
%                         largura fixa, para que colunas de diagrama fiquem
%                         alinhadas sem ajuste manual.
% crseta={cor}          : seta da familia musial na mesma cor semantica.
% passo={cor}           : circulo numerado com gradiente, para etapas.
% =============================================================================
\tikzset{
  crnode/.style n args={2}{musialBase, draw=#1!80!black,
                           top color=white, bottom color=#1!18,
                           text width=#2, align=center},
  crseta/.style={-{Latex[length=2.5mm,width=1.8mm]}, line width=0.8pt,
                 rounded corners=4pt, #1!85!black},
  passo/.style={circle, very thick, draw=#1!85!black,
                top color=white, bottom color=#1!30, text=#1!75!black,
                font=\scriptsize\bfseries\sffamily, inner sep=1.6pt,
                minimum size=5.2mm},
}

% =============================================================================
% \bandaimg : moldura clara com sombra sobre figura de fundo branco. Resolve a
% costura da figura com a marca d'agua do slide. Usar sempre que a ilustracao
% tiver fundo branco puro.
%   \bandaimg{opcoes de includegraphics}{arquivo}
% =============================================================================
\newcommand{\bandaimg}[2]{%
  \begin{tikzpicture}
    \node[rectangle, rounded corners=3pt, inner sep=3pt, fill=white,
          draw=campografite!25, line width=0.5pt,
          blur shadow={shadow blur steps=4, shadow xshift=0.8pt,
                       shadow yshift=-0.8pt}]
      {\includegraphics[keepaspectratio,#1]{#2}};
  \end{tikzpicture}%
}

% =============================================================================
% \reservaimg : espaco reservado para ilustracao ainda nao produzida. Compila
% como moldura vazia com a legenda do que vai entrar ali. Ao existir o arquivo,
% trocar a chamada por \bandaimg com o mesmo primeiro argumento.
%   \reservaimg{largura}{altura}{descricao curta}
% =============================================================================
\newcommand{\reservaimg}[3]{%
  \begin{tikzpicture}
    \node[rectangle, rounded corners=3pt, minimum width=#1, minimum height=#2,
          fill=campografite!4, draw=campografite!30, line width=0.5pt,
          dash pattern=on 2.6pt off 2.2pt, align=center,
          text width={\dimexpr #1 - 10pt\relax},
          font=\scriptsize\sffamily\color{campografite!85}]
      {\faFile\quad reserva de ilustra\c{c}\~ao\\[2pt]\scriptsize #3};
  \end{tikzpicture}%
}

% =============================================================================
% \imgcanvas : sistema de coordenadas normalizado sobre uma figura, de (0,0)
% no canto inferior esquerdo a (1,1) no canto superior direito, para acrescentar
% rotulos VETORIAIS sem tocar no arquivo rasterizado.
% =============================================================================
\newenvironment{imgcanvas}[2]{%
  \begin{tikzpicture}
    \node[anchor=south west, inner sep=0] (img) at (0,0)
      {\includegraphics[keepaspectratio,#1]{#2}};
    \begin{scope}[x={(img.south east)}, y={(img.north west)}]
}{%
    \end{scope}
  \end{tikzpicture}%
}
